% Preâmbulo compartilhado pelos relatórios T02 - Redes de Dados
\documentclass[11pt,a4paper,oneside]{article}

\usepackage[utf8]{inputenc}
\usepackage[T1]{fontenc}
\usepackage[brazilian]{babel}
\usepackage{lmodern}
\usepackage{microtype}

\usepackage[a4paper,margin=2.5cm,top=2.5cm,bottom=2.5cm]{geometry}
\usepackage{graphicx}
\usepackage{xcolor}
\usepackage{tabularx}
\usepackage{longtable}
\usepackage{booktabs}
\usepackage{array}
\usepackage{enumitem}
\usepackage{titlesec}
\usepackage{fancyhdr}
\usepackage{listings}
\usepackage{hyperref}
\usepackage{caption}
\usepackage{float}
\usepackage{tcolorbox}
\tcbuselibrary{skins,breakable}

% Cores
\definecolor{codebg}{RGB}{245,245,247}
\definecolor{codeborder}{RGB}{220,220,225}
\definecolor{codekey}{RGB}{0,90,160}
\definecolor{codecom}{RGB}{100,110,120}
\definecolor{codestr}{RGB}{40,140,80}
\definecolor{noteblue}{RGB}{30,80,160}
\definecolor{notebg}{RGB}{235,243,255}
\definecolor{warnyellow}{RGB}{170,120,0}
\definecolor{warnbg}{RGB}{255,248,225}

% Hyperref
\hypersetup{
  colorlinks=true,
  linkcolor=black,
  urlcolor=codekey,
  citecolor=black,
  pdfborder={0 0 0}
}

% Cabeçalho/rodapé
\pagestyle{fancy}
\fancyhf{}
\fancyhead[L]{\small T02 - Redes de Dados}
\fancyhead[R]{\small \leftmark}
\fancyfoot[C]{\thepage}
\renewcommand{\headrulewidth}{0.4pt}

% Listings (código)
\lstdefinestyle{shell}{
  backgroundcolor=\color{codebg},
  basicstyle=\ttfamily\footnotesize,
  keywordstyle=\color{codekey}\bfseries,
  commentstyle=\color{codecom}\itshape,
  stringstyle=\color{codestr},
  numbers=none,
  frame=single,
  rulecolor=\color{codeborder},
  framesep=4pt,
  xleftmargin=6pt,
  xrightmargin=6pt,
  breaklines=true,
  breakatwhitespace=false,
  showstringspaces=false,
  columns=fullflexible,
  keepspaces=true,
  literate=
    {á}{{\'a}}1 {é}{{\'e}}1 {í}{{\'i}}1 {ó}{{\'o}}1 {ú}{{\'u}}1
    {Á}{{\'A}}1 {É}{{\'E}}1 {Í}{{\'I}}1 {Ó}{{\'O}}1 {Ú}{{\'U}}1
    {ã}{{\~a}}1 {õ}{{\~o}}1 {Ã}{{\~A}}1 {Õ}{{\~O}}1
    {â}{{\^a}}1 {ê}{{\^e}}1 {ô}{{\^o}}1 {Â}{{\^A}}1 {Ê}{{\^E}}1 {Ô}{{\^O}}1
    {ç}{{\c{c}}}1 {Ç}{{\c{C}}}1
    {à}{{\`a}}1
}
\lstset{style=shell}

% Caixas de nota/observação/atenção
\newtcolorbox{notebox}{
  enhanced, breakable,
  colback=notebg, colframe=noteblue,
  boxrule=0.6pt, arc=2pt,
  left=8pt, right=8pt, top=4pt, bottom=4pt,
  fonttitle=\bfseries
}
\newtcolorbox{warnbox}{
  enhanced, breakable,
  colback=warnbg, colframe=warnyellow,
  boxrule=0.6pt, arc=2pt,
  left=8pt, right=8pt, top=4pt, bottom=4pt,
  fonttitle=\bfseries
}

% Imagem local
% Uso: \imgph{largura}{caminho}
\newcommand{\imgph}[2]{%
  \begin{figure}[H]\centering
    \includegraphics[width=#1]{#2}%
  \end{figure}
}

% Espaçamento de seções
\titleformat{\section}{\Large\bfseries\color{noteblue}}{\thesection.}{0.6em}{}
\titleformat{\subsection}{\large\bfseries}{\thesubsection}{0.6em}{}
\titleformat{\subsubsection}{\normalsize\bfseries}{\thesubsubsection}{0.6em}{}

% Capa
\newcommand{\capa}[2]{%
  \begin{titlepage}
    \centering
    \vspace*{1.5cm}
    % Logo: coloque o arquivo em .tex/img/logo.png
    \IfFileExists{imgs/logo.png}{\includegraphics[width=4cm]{imgs/logo.png}}{%
      \fbox{\parbox[c][3cm][c]{4cm}{\centering\footnotesize\itshape [logo.png]}}}
    \vfill
    {\large\bfseries T02 -- Redes de Dados\par}
    \vspace{1cm}
    {\Huge\bfseries\color{noteblue} #1\par}
    \vspace{0.8cm}
    {\large #2\par}
    \vfill
    {\large \today\par}
    \vspace{1cm}
  \end{titlepage}
}


\begin{document}

\capa{Relatório 10}{TLS em HTTP com Wireshark (VMs no VirtualBox)}

\tableofcontents
\newpage

\noindent
Este relatório tem como objetivo \textbf{montar um laboratório em NAT Network (rede NAT do VirtualBox)} para demonstrar o uso de \textbf{TLS} no protocolo \textbf{HTTP}, comparando o tráfego em texto claro (\textbf{HTTP sem TLS}) e criptografado (\textbf{HTTPS com TLS}) e \textbf{analisando as diferenças no Wireshark}.

\begin{itemize}
  \item \textbf{Material de apoio:} \url{https://github.com/vin1sss/T02---Redes-de-Dados-1/}
\end{itemize}

\section{Introdução}

Nesta atividade, você configurará duas VMs (\textbf{user 1 - Debian} e \textbf{user 2 - Debian}) em uma \textbf{NAT Network} do VirtualBox para observar, via Wireshark no \textbf{user 2 - Debian}, como o TLS protege o conteúdo das mensagens HTTP. Em texto claro, o payload é legível; com TLS, apenas metadados e o \textbf{handshake} ficam visíveis.

O TLS cria um canal cifrado sobre TCP por meio do \textit{handshake}, com negociação de chaves de sessão temporárias para proteger o tráfego de aplicação. Na prática, isso reduz a exposição a interceptação e leitura de conteúdo por terceiros no caminho.

\textbf{Pilares abordados:} foco em \textbf{Confidencialidade} (com menção a Autenticidade/Integridade via certificados e MAC do TLS).

\section{Conceitos e Fundamentos}

\begin{itemize}
  \item \textbf{HTTP x HTTPS:} HTTPS = HTTP sobre TLS (camada de segurança entre transporte e aplicação).
  \item \textbf{TLS (Handshake):} ClientHello $\rightarrow$ ServerHello (+ certificado) $\rightarrow$ \textit{key exchange} $\rightarrow$ chaves de sessão $\rightarrow$ comunicação criptografada.
  \item \textbf{Certificados:} neste lab usaremos certificado \textbf{autoassinado} no Apache.
  \item \textbf{Wireshark:} filtros úteis \texttt{http}, \texttt{tls}, \texttt{tcp.port == 443}, \texttt{tls.handshake}.
\end{itemize}

\section{Ambiente e Topologia}

\textbf{VirtualBox --- NAT Network:} \texttt{NatNetwork}

\begin{warnbox}
\textbf{Atenção:} neste relatório usamos \textbf{NAT Network} (rede compartilhada entre VMs). Isso é diferente de \textbf{NAT} simples por VM (\textit{Attached to: NAT}), que normalmente não permite comunicação direta entre as VMs.
\end{warnbox}

\subsection{Resumo do laboratório}

\begin{tabularx}{\linewidth}{lX}
\toprule
\textbf{Elemento} & \textbf{Definição} \\
\midrule
Rede do VirtualBox & \textbf{NAT Network} \texttt{NatNetwork} \\
user 1 - Debian & Servidor Apache, usando as portas \textbf{80} e \textbf{443} \\
user 2 - Debian & Cliente de testes com \texttt{curl} e \textbf{Wireshark} \\
Adaptadores das VMs & Uma placa por VM, ligada à \textbf{NAT Network} \\
Recursos das VMs & \textbf{6 GB de RAM} e \textbf{3 núcleos de CPU} \\
\bottomrule
\end{tabularx}

\subsection{Dados de referência usados no relatório}

\begin{tabularx}{\linewidth}{lX}
\toprule
\textbf{Item} & \textbf{Uso} \\
\midrule
Senha das máquinas virtuais & \texttt{pytest} \\
Senha solicitada pelo \texttt{sudo} & \texttt{pytest} \\
\texttt{<IP\_USER1>} & IP atual do \textbf{user 1 - Debian} \\
\texttt{<IP\_USER2>} & IP atual do \textbf{user 2 - Debian} \\
\bottomrule
\end{tabularx}

\subsection{Pacotes usados no laboratório}

\begin{itemize}
  \item \textbf{user 1 - Debian:} \texttt{apache2}, \texttt{openssl}
  \item \textbf{user 2 - Debian:} \texttt{curl}, \texttt{wireshark}
\end{itemize}

\subsection{Guia auxiliar de VirtualBox}

\begin{notebox}
Se não estiver conseguindo copiar e colar comandos entre o computador host e a VM, redimensionar a tela corretamente ou usar melhor a integração com o VirtualBox, consulte o guia \texttt{VirtualBox\_Guest\_Additions.md} para instalar o \textbf{VirtualBox Guest Additions}.
\end{notebox}

\section{Instalação e Preparação}

\begin{notebox}
Substitua as ocorrências de \texttt{<IP\_USER1>} e \texttt{<IP\_USER2>} pelos endereços obtidos nas suas VMs.
\end{notebox}

\subsection{Preparar a rede no VirtualBox}

\begin{enumerate}
\item \textbf{Abrir o gerenciador de redes NAT:}
  \begin{itemize}
    \item No VirtualBox, clique em \textbf{``Arquivos''}
    \item Clique em \textbf{``Ferramentas''}
    \item Clique em \textbf{``Gerenciador de Rede''}
    \item Clique em \textbf{``Redes NAT''}
  \end{itemize}

\item \textbf{Se já existir \texttt{NatNetwork}, apenas confira se está configurada:}
  \begin{itemize}
    \item \textbf{Name:} \texttt{NatNetwork}
    \item \textbf{IPv4 Prefix:} verifique o prefixo (ex.: \texttt{10.0.2.0/24})
    \item \textbf{Habilitar DHCP:} Marcado
  \end{itemize}

\item \textbf{Se NÃO existir \texttt{NatNetwork}: criar uma:}
  \begin{itemize}
    \item Clique em \textbf{``Criar''}
    \item \textbf{Name:} \texttt{NatNetwork}
    \item \textbf{IPv4 Prefix:} ex.: \texttt{10.0.2.0/24} (padrão)
    \item Marque \textbf{Habilitar DHCP}
    \item Confirme com \textbf{Aplicar}
  \end{itemize}

\imgph{0.6\linewidth}{imgs/a890ec72-c5da-49a0-81c3-e21ee7deb404.png}

\item \textbf{Vincular cada VM à \texttt{NatNetwork}:}
  \begin{itemize}
    \item Para \textbf{user 1 - Debian} e \textbf{user 2 - Debian} $\rightarrow$ \textbf{Configurações} $\rightarrow$ \textbf{Rede}
    \item \textbf{Adaptador 1} $\rightarrow$ \textbf{Conectado a:} \textit{NAT Network} $\rightarrow$ \textbf{Nome:} \texttt{NatNetwork}
    \item OK
  \end{itemize}

\imgph{0.6\linewidth}{imgs/3a4a1d8a-63b4-4c48-b06a-d65de72964a2.png}

\item \textbf{Ajustar recursos de hardware das VMs:}
  \begin{itemize}
    \item VirtualBox $\rightarrow$ \textbf{Configurações} $\rightarrow$ \textbf{Sistema} $\rightarrow$ \textbf{Placa-mãe} $\rightarrow$ \textbf{Memória Base: 6144 MB (6 GB)}
    \item VirtualBox $\rightarrow$ \textbf{Configurações} $\rightarrow$ \textbf{Sistema} $\rightarrow$ \textbf{Processador} $\rightarrow$ \textbf{Processadores: 3}
    \item Após ajustar os recursos, \textbf{iniciar} as VMs
  \end{itemize}

\imgph{0.4\linewidth}{imgs/4e1de125-9012-4b75-b395-e3aa73a179a8.png}
\end{enumerate}

\begin{notebox}
\textbf{Ao final desta etapa:} as duas VMs devem estar ligadas, vinculadas à \texttt{NatNetwork} e prontas para receber IP via DHCP.
\end{notebox}

\subsection{Conferir IPs e validar a comunicação}

\begin{enumerate}
\item Em \textbf{cada VM Debian}, execute:
\begin{lstlisting}
ip a
\end{lstlisting}

\item Identifique o endereço IPv4 usado na interface conectada à \texttt{NatNetwork} e verifique se os IPs de \textbf{user 1 - Debian} e \textbf{user 2 - Debian} pertencem à \textbf{mesma faixa de rede}.
  \begin{itemize}
    \item Exemplo: se a \texttt{NatNetwork} estiver em \texttt{10.0.2.0/24}, os dois IPs devem seguir \texttt{10.0.2.x/24}.
    \item Anote o IP do servidor como \texttt{<IP\_USER1>} e o do cliente como \texttt{<IP\_USER2>}.
  \end{itemize}

\imgph{0.8\linewidth}{imgs/70c2e67a-184a-487b-a232-36a1759252e1.png}

\item \textbf{Se os IPs estiverem compatíveis}, avance para o teste de comunicação.

\item \textbf{Se os IPs não estiverem na mesma rede}, execute em \textbf{cada VM Debian}:
\begin{lstlisting}
IFACE=$(ip route | awk '/default/ {print $5; exit}')
sudo dhclient -r "$IFACE" || true
sudo ip addr flush dev "$IFACE"
sudo dhclient "$IFACE"
ip -4 a show "$IFACE"
\end{lstlisting}

\begin{notebox}
A interface padrão em Debian no VirtualBox é normalmente \texttt{enp0s3}. Confirme com \texttt{ip route} ou \texttt{ip a}.
\end{notebox}

\item \textbf{Testar comunicação entre as VMs com \texttt{ping}:}

No \textbf{user 2 - Debian}, execute:
\begin{lstlisting}
ping -c2 <IP_USER1>
\end{lstlisting}

No \textbf{user 1 - Debian}, execute:
\begin{lstlisting}
ping -c2 <IP_USER2>
\end{lstlisting}

\imgph{0.8\linewidth}{imgs/e537c6a0-607b-4553-abec-e13e0b7375ad.png}
\end{enumerate}

\begin{notebox}
\textbf{Ao final desta etapa:} as duas VMs devem estar na mesma rede e responder ao \texttt{ping} entre si.
\end{notebox}

\subsection{Preparar o servidor --- user 1 - Debian}

\subsubsection*{A) Instalar dependências}

\begin{lstlisting}
sudo apt update && sudo apt install -y apache2 openssl
\end{lstlisting}

\imgph{0.55\linewidth}{imgs/1592bc42-1a35-46b9-8e64-ac8779790f87.png}

\subsubsection*{B) Criar página de teste com conteúdo identificável}

\begin{lstlisting}
sudo mkdir -p /var/www/html
echo "<h1>HELLO_TLS_HTTP</h1>" | sudo tee /var/www/html/index.html
\end{lstlisting}

\imgph{0.5\linewidth}{imgs/1d5ca830-2b83-4197-9bb3-a7fd1f6852c1.png}

\subsubsection*{C) Ativar o Apache e verificar a porta HTTP}

\begin{lstlisting}
sudo systemctl enable --now apache2
ss -tulpn | grep :80
\end{lstlisting}

\begin{notebox}
\textbf{Resultado esperado:} deve aparecer uma linha indicando que o Apache está em estado \textbf{LISTEN} na porta \textbf{80}.
\end{notebox}

\imgph{0.55\linewidth}{imgs/2c5dc094-4ede-4deb-90fd-0c3007cca55d.png}

\begin{notebox}
\textbf{Ao final desta etapa:} o servidor deve estar respondendo em \textbf{HTTP} na porta \textbf{80}.
\end{notebox}

\subsection{Preparar o cliente --- user 2 - Debian}

\subsubsection*{A) Instalar \texttt{curl} e Wireshark}

\begin{notebox}
Observação: escolha ``Yes'' na configuração do Wireshark.
\end{notebox}

\begin{lstlisting}
sudo apt update
sudo apt install -y curl wireshark
\end{lstlisting}

\imgph{0.5\linewidth}{imgs/92dbd981-4482-47f4-a52c-a111d0a10500.png}

\subsubsection*{B) Abrir o Wireshark e identificar a interface de captura}

\begin{lstlisting}
sudo wireshark
\end{lstlisting}

Esse comando \textbf{apenas abre o programa}. Na tela inicial, identifique a \textbf{interface da NAT Network} (a que mostra o IP do \textbf{user 2 - Debian}, costuma ser \texttt{enp0s3}) e deixe o Wireshark pronto para iniciar a captura no próximo passo.

\imgph{0.7\linewidth}{imgs/9e375ed7-f018-4bb6-9ffd-9ccbb67e950a.png}

Durante a análise dos pacotes, use como referência o filtro:
\begin{lstlisting}
ip.addr == <IP_USER1>
\end{lstlisting}

\begin{notebox}
\textbf{Ao final desta etapa:} o cliente deve estar com \texttt{curl} instalado, o Wireshark aberto e a interface correta já identificada.

\textbf{Observação:} os filtros indicados neste relatório são \textbf{filtros de exibição} no Wireshark, usados para localizar melhor os pacotes já capturados.
\end{notebox}

\section{Procedimentos (Passo a Passo)}

\subsection*{Checklist antes dos testes}

Antes de iniciar os cenários, confirme:
\begin{itemize}
  \item o IP do servidor está anotado como \texttt{<IP\_USER1>};
  \item o Apache responde em \textbf{HTTP} na porta \textbf{80};
  \item o Wireshark está aberto no \textbf{user 2 - Debian};
  \item a interface da \textbf{NAT Network} já foi identificada (costuma ser \texttt{enp0s3}).
\end{itemize}

\subsection{Cenário 1 --- HTTP sem TLS (porta 80)}

\subsubsection*{A) Preparar a captura}

\textbf{No user 2 - Debian:} com o Wireshark já aberto, \textbf{inicie a captura} na interface da \textbf{NAT Network} identificada anteriormente.

\imgph{0.3\linewidth}{imgs/480f3519-fe3c-4dfd-8cec-329fc63a8c0d.png}

\imgph{0.6\linewidth}{imgs/e8be4387-047f-4705-a085-4eafb714164f.png}

\begin{notebox}
O quadrado vermelho indica que a captura está ativa.
\end{notebox}

\subsubsection*{B) Executar o teste HTTP}

No \textbf{user 2 - Debian}, execute:
\begin{lstlisting}
curl -v http://<IP_USER1>
\end{lstlisting}

\begin{notebox}
\textbf{Resultado esperado:} o \texttt{curl} deve completar a requisição e exibir a resposta do Apache, incluindo \texttt{HELLO\_TLS\_HTTP}.
\end{notebox}

\imgph{0.45\linewidth}{imgs/3c0016d3-3879-4451-94a3-f9b2ad3f8a92.png}

\subsubsection*{C) Observar no Wireshark}

\begin{enumerate}
\item \textbf{Pause a captura} no Wireshark.

\imgph{0.2\linewidth}{imgs/cfedfeba-6545-4297-aa1b-530b14688ed8.png}

\item Aplique o filtro:
\begin{lstlisting}
http && ip.addr == <IP_USER1>
\end{lstlisting}

\item \textbf{O que observar:} requisição \texttt{GET / HTTP/1.1} e resposta \texttt{200 OK} com \textbf{payload legível} (HTML contendo \texttt{HELLO\_TLS\_HTTP}).

\imgph{0.7\linewidth}{imgs/43497e00-b9d0-436b-8072-859dcaa368e0.png}
\imgph{0.5\linewidth}{imgs/2310e98c-7f1a-40fc-887a-592cdceb22a9.png}
\end{enumerate}

\begin{notebox}
\textbf{Ao final deste cenário:} o tráfego HTTP deve aparecer em texto claro no Wireshark.
\end{notebox}

\subsection{Cenário 2 --- HTTP com TLS (HTTPS na porta 443)}

\subsubsection*{A) Preparar HTTPS no servidor --- user 1 - Debian}

\textbf{Ativar SSL no Apache:}
\begin{lstlisting}
sudo a2enmod ssl
sudo a2ensite default-ssl
sudo systemctl reload apache2
\end{lstlisting}

\imgph{0.5\linewidth}{imgs/46ea3fc5-5434-44a8-8773-2cb63ed55ee3.png}

\begin{notebox}
Caso apareça um erro relacionado ao Apache, ignore --- ele apenas indica que o Apache não estava rodando anteriormente.
\end{notebox}

\textbf{Gerar certificado autoassinado (1 ano):}
\begin{lstlisting}
sudo openssl req -x509 -nodes -days 365 -newkey rsa:2048 \
  -keyout /etc/ssl/private/apache.key \
  -out /etc/ssl/certs/apache.crt \
  -subj "/C=BR/ST=MG/L=Inatel/O=Lab/OU=NetSec/CN=<IP_USER1>"
\end{lstlisting}

\imgph{0.7\linewidth}{imgs/fa045612-e0ef-4120-9123-b8f47b2144ff.png}

\textbf{Apontar o vhost SSL para o certificado/chave e reiniciar Apache:}
\begin{lstlisting}
sudo sed -i 's#SSLCertificateFile .*#SSLCertificateFile /etc/ssl/certs/apache.crt#' /etc/apache2/sites-available/default-ssl.conf
sudo sed -i 's#SSLCertificateKeyFile .*#SSLCertificateKeyFile /etc/ssl/private/apache.key#' /etc/apache2/sites-available/default-ssl.conf
sudo systemctl restart apache2
ss -tulpn | grep :443
\end{lstlisting}

\begin{notebox}
\textbf{Resultado esperado:} Apache em estado \textbf{LISTEN} na porta \textbf{443}.
\end{notebox}

\imgph{0.7\linewidth}{imgs/baaaf721-a9ce-4249-8f87-748c2494a942.png}

\subsubsection*{B) Reiniciar a captura e executar o teste HTTPS --- user 2 - Debian}

No Wireshark, \textbf{reinicie a captura} na interface da \textbf{NAT Network} para registrar somente o tráfego do novo teste.

\imgph{0.6\linewidth}{imgs/069ce333-c2a3-46b4-8ada-6343845d21f3.png}
\imgph{0.6\linewidth}{imgs/e8be4387-047f-4705-a085-4eafb714164f.png}

Execute:
\begin{lstlisting}
curl -vk https://<IP_USER1>
\end{lstlisting}

\begin{notebox}
O \texttt{-k} é intencional no laboratório: ele permite seguir com certificado autoassinado. \textbf{Resultado esperado:} a conexão HTTPS deve funcionar mesmo com o certificado autoassinado.
\end{notebox}

\imgph{0.55\linewidth}{imgs/b612937e-31f3-40ec-9cfc-715101247637.png}

\subsubsection*{C) Observar no Wireshark}

\begin{enumerate}
\item \textbf{Pause a captura} no Wireshark.
\item Remova o filtro anterior (se houver).
\item Aplique o filtro:
\begin{lstlisting}
tls && ip.addr == <IP_USER1>
\end{lstlisting}
\item \textbf{O que observar:} pacotes de \textbf{handshake TLS} (ClientHello/ServerHello/Certificado) e \textbf{payload cifrado}.
\end{enumerate}

\imgph{0.8\linewidth}{imgs/033753e3-daae-4dfd-966a-ce779528bdf9.png}
\imgph{0.8\linewidth}{imgs/304bdb65-de32-46b1-a97e-19e5b85d3833.png}

\begin{notebox}
\textbf{Ao final deste cenário:} o conteúdo da aplicação não deve aparecer legível no Wireshark; apenas o handshake e os dados criptografados ficam visíveis.
\end{notebox}

\section{Verificação e Resultados}

\subsection{Checklist de sucesso do experimento}

\begin{itemize}
  \item \textbf{Cenário HTTP (porta 80):} \texttt{curl -v http://<IP\_USER1>} retorna conteúdo e o Wireshark mostra \texttt{GET}/\texttt{200 OK} com payload legível.
  \item \textbf{Cenário HTTPS (porta 443):} \texttt{curl -vk https://<IP\_USER1>} funciona, e no Wireshark aparecem \texttt{ClientHello}/\texttt{ServerHello} com payload cifrado.
\end{itemize}

\subsection{Quadro comparativo}

\begin{tabularx}{\linewidth}{lccXX}
\toprule
\textbf{Protocolo} & \textbf{Porta} & \textbf{TLS} & \textbf{Visível no Wireshark} & \textbf{Payload} \\
\midrule
HTTP  & 80  & Não & Métodos/URLs, cabeçalhos, corpo & \textbf{Legível} (\texttt{HELLO\_TLS\_HTTP}) \\
HTTPS & 443 & Sim & Handshake (ClientHello/ServerHello/Cert), \textit{Application Data} & \textbf{Ilegível} \\
\bottomrule
\end{tabularx}

\subsection{Filtros práticos (copiar/colar)}

\begin{itemize}
  \item HTTP claro: \texttt{http}
  \item HTTPS: \texttt{tls} (ou \texttt{tcp.port == 443})
  \item Handshake específico: \texttt{tls.handshake}
  \item Verificar string no tráfego claro: \texttt{frame contains "HELLO\_TLS\_HTTP"}
\end{itemize}

\section{Conclusão}

Foi demonstrado, de forma prática, que a adoção de \textbf{TLS} em HTTP protege a \textbf{confidencialidade} das mensagens. Enquanto no tráfego sem TLS foi possível ler o payload em claro, no tráfego com TLS o Wireshark exibiu apenas metadados do handshake e pacotes cifrados.

\section*{Apêndice --- Troubleshooting rápido}
\addcontentsline{toc}{section}{Apêndice --- Troubleshooting rápido}

\begin{itemize}
  \item \textbf{Não vejo tráfego no Wireshark:} selecione a interface da NAT Network correta e gere tráfego com \texttt{curl} no user 2 - Debian.
  \item \textbf{HTTPS não sobe:} confirme \texttt{a2enmod ssl}, \texttt{a2ensite default-ssl}, caminhos de certificado no \texttt{default-ssl.conf} e reinício do Apache.
  \item \textbf{\texttt{curl -vk https://<IP\_USER1>} falha:} verifique se a porta 443 está em LISTEN (\texttt{ss -lntp | grep :443}).
  \item \textbf{\texttt{ping} entre as VMs falha:} confirme que ambas estão em \textbf{NAT Network (\texttt{NatNetwork})}, mesma faixa de rede e endereço correto de destino.
  \item \textbf{VMs não se enxergam:} confirme ambas em \textbf{NAT Network (\texttt{NatNetwork})} com DHCP ON.
  \item \textbf{Wireshark não abre ou não captura corretamente:} execute novamente com \texttt{sudo wireshark} e confirme que a interface da NAT Network foi selecionada.
\end{itemize}

\end{document}
